\chapter{Ferramenta de Validação em Produção}

\section{O que é a Ferramenta de Validação em Produção (FVP)?}
A FVP é a plataforma utilizada para a realização de testes nas implementações em produção das instituições participantes e dos fornecedores de testes, com o objetivo de verificar a conformidade dos produtos e APIs do Open Finance em relação aos padrões técnicos definidos.
\medskip
A plataforma foi criada e é gerenciada pela empresa Raidiam, sob demanda da Associação Open Finance Brasil. 

\subsection{Como funciona a ferramenta?}
A FVP realiza testes automatizados e manuais em ambientes em produção, verificando se os fluxos e as respostas das APIs estão alinhados com as especificações técnicas do Open Finance Brasil. Isso inclui desde testes básicos de registro e atualização de clientes (DCR/DCM) até jornadas completas de compartilhamento de dados e iniciação de pagamentos, sempre de forma controlada e segura. 
\medskip
As informações de execução são disponibilizadas nos logs dos testes e são automaticamente removidas da ferramenta após 1 hora do início da execução. Os logs exportados podem conter dados sensíveis, sendo responsabilidade do usuário garantir o correto manuseio e compartilhamento desses arquivos, para evitar possíveis vazamentos. 

\section{Modalidades da FVP Manual}

A FVP Manual opera em duas modalidades distintas: \textbf{Testes Abertos}, destinados às instituições participantes, e \textbf{Testes Restritos}, exclusivos para fornecedores credenciados.

\begin{figure}[H]
\centering % para centralizarmos a figura
\figframe{width=0.8\textwidth}{fvpimage2.png}
\label{fig:fvpimage2}
\end{figure}

\begin{itemize}
    \item \textbf{Testes Abertos} -  Executados pelas instituições participantes.
    
    \item \textbf{Testes Restritos} - Executados por fornecedores contratados pela Associação Open Finance.
\end{itemize}

\subsection{FVP Manual - Testes Abertos}

Os Testes Abertos permitem que as instituições participantes executem testes em ambiente de produção em seus próprios servidores, não sendo necessária a atuação de fornecedores contratados pela Associação para a sua execução.
\medskip
Ao final de cada ciclo, será emitida notificações de irregularidade para as instituições que não executarem ou não obtiverem êxito nos Testes Abertos previstos.
As notificações de irregularidade somente podem ser encerradas após a execução bem-sucedida dos testes, conforme validação realizada pela Associação.
\medskip
Os testes exigidos com sucesso para cada produto sempre serão disponibilizados na area do desenvolvedor - \href{https://openfinancebrasil.atlassian.net/wiki/spaces/OF/pages/553255015/FVP+Manual#2---FVP-Manual---Testes-Abertos}{FVP Manual - Testes Abertos}

\subsection{FVP Manual - Testes Restritos}

Quando o fornecedor contratado pela estrutura possuir conta aberta na Instituição ou Marca, o fornecedor será o responsável pela execução dos testes da FVP Manual.
\medskip
Caso seja identificada qualquer falha durante a execução nos servidores das instituições, será aberto um ticket no Service Desk, contendo as evidências técnicas da falha e direcionado à instituição responsável.
\medskip
O encerramento do ticket está condicionado à reexecução bem-sucedida do módulo de teste da FVP, a qual deve ser realizada exclusivamente pelo fornecedor responsável pela execução dos testes.
\medskip
Testes realizados pela instituição no contexto onde fornecedores possuem contas abertas não serão aceitos como evidência de sucesso. 
\medskip
Os testes exigidos com sucesso para cada produto sempre serão disponibilizados na area do desenvolvedor - \href{https://openfinancebrasil.atlassian.net/wiki/spaces/OF/pages/553255015/FVP+Manual#3---FVP-Manual---Testes-Restritos}{FVP Manual - Testes Restritos}

\section{Modelos de Execução}

A FVP adota dois modelos distintos para a execução de testes, definidos de acordo com a duração e a continuidade do processo. O modelo Imediato é voltado a execuções que se completam no mesmo dia, enquanto o modelo Agendado ou de Longa Duração permite que a execução se estenda por vários dias de forma automatizada.

\begin{itemize}
    \item \textbf{Testes Imediatos} - São concluídas em uma única rodada, finalizando no mesmo dia em que são iniciadas.
    
    \item \textbf{Testes Agendados/Longa Duração} - Permitem iniciar o fluxo em um dia e continuar automaticamente nos dias subsequentes.
\end{itemize}

\begin{figure}[H]
    \centering
    \figframe{width=0.8\textwidth}{imediatoxlonga.png}
    \label{fig:imediatoxlonga}
\end{figure}

